\begin{abstract}
% --- SOFT DELETE: Original abstract (Batch 1) ---
% The high cost of proprietary Big Data infrastructure and software licenses presents a significant 
% barrier to entry for academic research and small organisations. This paper proposes a cost-effective, 
% open-source Data Lake solution deployed entirely on the Oracle Cloud Infrastructure (OCI) Free Tier. 
% The methodology leverages a cluster of ARM compute instances, autonomously orchestrated by Apache 
% Ambari within the Open Source Data Platform (ODP) ecosystem. The proposed architecture establishes a 
% comprehensive data pipeline, integrating ingestion via Apache NiFi and Kafka, distributed storage 
% with HDFS, and computational analysis using Apache Spark and Hive. To validate the 
% viability and resilience of the provisioned environment, stress-testing was conducted 
% utilising the Star Schema Benchmark (SSB) catalog. Analytical evaluations compared traditional 
% disk-bound batch processing in Hive against distributed execution in Spark. Performance metrics 
% revealed that, despite infrastructural constraints, the Spark engine consistently yielded 
% a performance edge, frequently outperforming Hive operations by margins exceeding 200 
% seconds per query. The empirical results validate the feasibility of this constraint-driven, 
% Infrastructure-as-Code (IaC) model as a fully reproducible and functional multi-node distributed 
% system, democratising access to Big Data technologies for research and educational 
% purposes.
% --- END SOFT DELETE ---
% --- SOFT DELETE: Revised abstract v1 (Batch 5) ---
% [see git diff for previous version; trimmed to <150 words and aligned caching
% claim with actual benchmark findings]
% --- END SOFT DELETE ---
The high cost of Big Data infrastructure limits access for academic and smaller
organisations. This paper presents an open-source Infrastructure-as-Code
framework that automates the deployment of a Hadoop-based Data Lake on Oracle
Cloud Infrastructure Free Tier ARM instances. Using Terraform and Ansible, the
framework provisions a four-node ARM cluster and installs the OpenSource Data
Platform stack via Apache Ambari, supporting multiple installation profiles. The
Star Schema Benchmark (SF\,=\,10) was executed under three
configurations---Hive, Spark, and Spark with caching---to evaluate query
performance on constrained hardware. Spark outperformed Hive by mean margins of
282 to 2\,112~seconds per query. In-memory caching degraded performance for
simpler queries owing to insufficient memory for the largest fact table, and
yielded improvements only for complex multi-join workloads. These results confirm
the viability of zero-cost, reproducible distributed analytics for research and
education.
\end{abstract}

% --- SOFT DELETE: Original keywords (Batch 1) ---
% \KEYWORD{Big Data; Data Lake; Infrastructure as Code; Apache Hadoop; Apache Spark; Cloud Computing; ARM Architecture; Distributed Systems; Star Schema Benchmark; Cost Optimisation.}
% --- END SOFT DELETE ---

\KEYWORD{Big Data, Data Lake, Infrastructure as Code, Apache Hadoop, Apache Spark, Cloud Computing, ARM Architecture, Distributed Systems, Star Schema Benchmark, Cost Optimisation, Free Tier, Oracle Cloud Infrastructure, Apache Ambari, Reproducibility.}
